%!TEX program = lualatex
\documentclass[12pt]{article}

\usepackage[portuguese]{babel}
\usepackage{fontspec}
\setmainfont{Latin Modern Roman}
\usepackage[a4paper,margin=2.5cm]{geometry}
\usepackage{setspace}

\begin{document}

\section*{Roteiro de apresentação — Pessoa 3}

\textbf{Discussão e resultados III}

Uma possível abordagem de mitigação envolve a construção de estruturas de proteção costeira, como quebra-mares.

Essas estruturas atuam reduzindo a energia das ondas incidentes na faixa litorânea.

\vspace{0.5em}

\textbf{Discussão e resultados IV}

Entre as soluções possíveis, destaca-se o uso de unidades artificiais, como os elementos Core-Loc.

Essas estruturas aumentam a estabilidade hidráulica e promovem a dissipação da energia das ondas.

No entanto, sua eficácia depende das condições locais e de estudos técnicos específicos.

\vspace{0.5em}

\textbf{Discussão e resultados V}

A implantação dessas estruturas exige controle geométrico rigoroso, com uso de técnicas de georreferenciamento.

O acompanhamento por meio de levantamentos “as built” permite avaliar o desempenho da obra ao longo do tempo.

\vspace{0.5em}

\textbf{Conclusões I}

Casos de sucesso mostram que estruturas costeiras podem reduzir a erosão quando bem dimensionadas.

Projetos compatíveis com a dinâmica local têm maior probabilidade de sucesso.

\vspace{0.5em}

\textbf{Conclusões II}

Diferentes configurações geométricas podem ser utilizadas, como espigões e quebra-mares.

A escolha da solução depende das características específicas do sistema costeiro.

\vspace{0.5em}

\textbf{Conclusões III}

A resposta da linha de costa às estruturas pode ser interpretada a partir do conceito de equilíbrio morfodinâmico.

A linha tracejada representa a posição projetada de equilíbrio da linha de costa após a alteração do campo de ondas.

\vspace{0.5em}

\textbf{Encerramento}

Em síntese, a erosão em Atafona resulta de um desequilíbrio hidrossedimentar.

A mitigação do problema exige uma abordagem integrada, considerando processos fluviais, marinhos e antrópicos.

Obrigado.

\end{document}